\section{Canonical module linearization of actual invariant Cartier ideals}
\begin{lemma}[Whole pullback inclusion and actual composition]
\label{lem:cartier-whole-pullback-inclusion}
\uses{lem:actual-ideal-module-canonical-pullback-comp,thm:actual-cartier-module-global-flat-pullback}
\leanok
\lean{AlgebraicGeometry.Scheme.IdealSheafData.pullbackModuleInclusion,AlgebraicGeometry.Scheme.IdealSheafData.pullbackModuleInclusion_eq,AlgebraicGeometry.Scheme.IdealSheafData.pullbackModuleInclusion_comp,AlgebraicGeometry.Scheme.IdealSheafData.IsInvertible.pullbackModuleInclusion_mono}
The actual pullback of the original ideal inclusion followed by the genuine
structure-module comparison is its canonical geometric comap inclusion.
It respects native whole pullback composition. For a Cartier ideal under
a flat morphism this whole module inclusion is a monomorphism.
\end{lemma}
\begin{proof}
Use the actual canonical ideal-module inclusion triangle for the first
identity. Naturality of native pullback composition moves the original
ideal inclusion past that composition isomorphism, and the genuine
structure-module unit composition identity gives the second assertion.
For the last assertion the actual canonical Cartier comparison is an
isomorphism by global flat base change, and the geometric ideal inclusion
is an actual monomorphism. Their composite is therefore a monomorphism.
\end{proof}

\begin{lemma}[Canonical comparison of two whole invariant Cartier pullbacks]
\label{lem:cartier-invariant-module-iso}
\uses{lem:cartier-whole-pullback-inclusion,lem:actual-ideal-module-equality-iso}
\leanok
\lean{AlgebraicGeometry.Scheme.IdealSheafData.IsInvertible.invariantModuleIso,AlgebraicGeometry.Scheme.IdealSheafData.IsInvertible.invariantModuleIso_hom_comp_canonical,AlgebraicGeometry.Scheme.IdealSheafData.IsInvertible.invariantModuleIso_hom_inclusion,AlgebraicGeometry.Scheme.IdealSheafData.IsInvertible.invariantModuleIso_hom_unique}
If two actual flat scheme morphisms from the same original source pull
back the entire Cartier ideal to the same geometric ideal, their actual
module pullbacks have a canonical isomorphism. It preserves the entire
included structure-module section map and is uniquely determined by that
whole inclusion identity.
\end{lemma}
\begin{proof}
Compose the genuine canonical flat Cartier comparison for the first
morphism, the module isomorphism from the literal entire ideal equality,
and the inverse of the genuine canonical comparison for the second
morphism. Cancelling its actual inverse proves the stated canonical-map
formula. Compose with the actual final ideal inclusion; whole ideal
equality preserves the original inclusion, giving the true ambient
inclusion identity. The second whole pullback inclusion is a monomorphism
by the preceding lemma, so cancellation proves uniqueness. No module
comparison or linearization is assumed.
\end{proof}

\begin{lemma}[Transitivity and actual pullback naturality]
\label{lem:cartier-invariant-module-coherence}
\uses{lem:cartier-invariant-module-iso,lem:cartier-whole-pullback-inclusion}
\leanok
\lean{AlgebraicGeometry.Scheme.IdealSheafData.IsInvertible.invariantModuleIso_refl,AlgebraicGeometry.Scheme.IdealSheafData.IsInvertible.invariantModuleIso_trans,AlgebraicGeometry.Scheme.IdealSheafData.IsInvertible.invariantModuleIso_pullback,AlgebraicGeometry.Scheme.IdealSheafData.IsInvertible.invariantModuleIso_congr,AlgebraicGeometry.Scheme.IdealSheafData.pullbackComparisonIso,AlgebraicGeometry.Scheme.IdealSheafData.IsInvertible.invariantModuleIso_pullback_iso}
These canonical invariant Cartier module isomorphisms have the identity
and transitivity laws. Pulling one back along any actual morphism whose two composites are flat
agrees with the canonical comparison of the composite maps, conjugated
by the genuine native pullback-composition isomorphisms.
\end{lemma}
\begin{proof}
For identity and transitivity, compose the actual candidate map with the
whole final pullback inclusion. The preceding normalization and actual
associativity reduce it to the original first inclusion. Uniqueness
therefore gives the claimed genuine entire sheaf map equality. For
pullback, postcompose both candidates with native composition for the
second morphism and its whole pullback inclusion. This is a monomorphism.
The true pullback-inclusion composition identity turns the pulled-back
side into the actual pullback of its normalization followed by the
original structure-module comparison. Functoriality and the actual
normalization identify it with the whole first composite inclusion.
The other side gives that same map by normalization for the genuine
composite ideal equality. Cancel the actual final monomorphism. This is
whole scheme functoriality, before specialization to the actual group
unit and multiplication diagrams. In particular the base morphism itself
need not be flat, as is needed for the actual group unit section. When
the two original scheme maps are literally equal, equality induction and
the true identity law give exactly the native pullback equality
isomorphism, with no chosen comparison map.
\end{proof}
